\section{Fermat型の普遍混合中心三次公式}
\label{sec:pg-fermat-mixed-universal}
本節では代数的な普遍中心元の公式を直接構成するため、
Fermat型の三変数Jacobianブロックの\emph{全混合係数}に対する
有限多項式の中心三次元を構成する。
この式は論文Bに関連する既存技術稿に書かれていたが、
これまでの独立検証は有限体特殊化に留まっていた。
ここでは六つの係数を互いに独立な記号として保ち、
$\hbar$も不定元とした全係数恒等式の証明を与える。

\begin{theorem}[普遍Fermat混合中心三次元]
\label{thm:pg-fermat-mixed-universal}
$R$を標数零の可換次数付き$\CC$-代数とし、
$A,B,C\in R_1$、$D,E,H\in R_2$、$G\in R_3$を任意に取る。
\begin{equation}
 F=\frac{x^3+y^3+z^3}{3}
   +Axy+Byz+Czx+Dx+Ey+Hz+G
 \label{eq:pg-fermat-mixed-F}
\end{equation}
とおき、
\begin{equation}
 [x,y]=\hbar(z^2+By+Cx+H),\quad
 [y,z]=\hbar(x^2+Ay+Cz+D),\quad
 [z,x]=\hbar(y^2+Ax+Bz+E)
 \label{eq:pg-fermat-mixed-rel}
\end{equation}
で定まる$R[\hbar]$-代数を$\mathcal A_F$とする。
$T=xyz+xzy+yxz+yzx+zxy+zyx$、
$\sym(xy)=(xy+yx)/2$などと書く。
このとき
\begin{align}
\Omega_{F,\hbar}={}&
 \frac{x^3+y^3+z^3}{3}-\frac{\hbar^2}{9}T
 +A\sym(xy)+B\sym(yz)+C\sym(zx)\notag\\
 &+\left(D+\frac{\hbar^2B^2}{6}-\frac{\hbar^4AC}{18}\right)x
 +\left(E+\frac{\hbar^2C^2}{6}-\frac{\hbar^4AB}{18}\right)y\notag\\
 &+\left(H+\frac{\hbar^2A^2}{6}-\frac{\hbar^4BC}{18}\right)z
 +\frac{\hbar^4}{18}(Az^2+Bx^2+Cy^2)+G
 \label{eq:pg-fermat-mixed-Omega}
\end{align}
は、少なくとも$\hbar=0$を含む共通のPBW開部分上で
$\mathcal A_F$の\emph{厳密な中心元}となり、
$\Omega_{F,0}=F$を満たす。
これは有限次ジェットの主張ではなく、全ての交換子が
厳密に零であるという代数的恒等式である。
特に式は$R$や中心変数の個数によらず有効である。
\end{theorem}
\begin{proof}
証明は有限次元の非可換自由語の\emph{普遍的な整数・有理関数恒等式}
で行う。係数は最初から独立な中心不定元としてよいので、
まず
\begin{equation}
 R_0=\mathbb Q[A,B,C,D,E,H],\qquad
 \deg(A,B,C)=1,\quad\deg(D,E,H)=2
\end{equation}
を考える。$x,y,z$も次数一、$\hbar$は次数零とする。
自由な語の各中央係数は可換に並べ、active変数の語順だけを保持する。

全次数四の自由語は、中央係数の重み$0,1,2,3,4$と
active語の長さ$4,3,2,1,0$ごとに列挙すると339個ある。
三つの関係式\eqref{eq:pg-fermat-mixed-rel}を
左右から全次数合計二の中央係数・active語で掛けると
162個の関係ベクトルを得る。
ここでactive語を長さ$a,b$に分けて
中央重み$c$を選び、$a+b+c=2$とする。
\begin{equation}
 \mathcal V_4=\mathrm{span}_{\mathbb Q(\hbar)}
  \{\text{全次数四の339語}\},\quad
 I_4=\mathrm{span}_{\mathbb Q(\hbar)}
  \{\text{上記162関係ベクトル}\}.
\end{equation}
これらは全次数四で必要な\emph{全ての}二次関係式の倍数である。

有理関数体$\mathbb Q(\hbar)$上で
160個余りの有限行列を正確に行簡約すると
$\dim I_4=144$であり、
\begin{equation}
 [\Omega_{F,\hbar},x],\quad
 [\Omega_{F,\hbar},y],\quad
 [\Omega_{F,\hbar},z]
 \quad\in I_4
 \label{eq:pg-fermat-mixed-certificate}
\end{equation}
が三つとも成立する。
全次数・全中央係数単項式・全active語を保持する
独立な計算実装
\texttt{computations/verify\_universal\_fermat\_mixed\_center.py}
は$\mathbb Q(\hbar)$の分数体演算だけを用いる。
乱数、有限体への特殊化、混合係数への代入は一切行わない。
スクリプトは339個の列、162個の関係ベクトル、
階数144と三つの零余剰を全て検査する。
したがってこれは任意の$A,B,C,D,E,H$への特殊化に先立つ
\emph{普遍的な恒等式}である。

分数体上での所属から$\hbar=0$を含む開部分へ
戻る操作にも、同じ有限次元行列で証明を与える。
$\hbar=0$では三関係式が通常の可換化関係式となり、
次数四の商は
$\mathbb Q[A,B,C,D,E,H,x,y,z]_4$である
（中央変数の重みは上記のとおり）。
中央重み$0,1,2,3,4$の単項式数は順に
$1,3,9,19,39$であり、
従って商の次元は
\begin{equation}
 1\cdot15+3\cdot10+9\cdot6+19\cdot3+39\cdot1
 =195,\qquad
 \operatorname{rank}I_4|_{\hbar=0}=339-195=144.
\end{equation}
この階数144は$\mathbb Q(\hbar)$上の階数144と等しく、
付属スクリプトでは双方を\emph{別々の行簡約}として検算する。
ゆえに離散付値環$\mathbb Q[\hbar]_{(\hbar)}$上の
この有限行列をSmith標準形にすると、
144個の非零不変因子は全て原点で単元である。
次数四の商は$\hbar$-torsion-freeとなり、
$I_4$は原点で$\hbar$-飽和している。
従って\eqref{eq:pg-fermat-mixed-certificate}は
$\hbar=0$で正則な係数の\emph{多項式関係式の恒等式}として成立する。
その後の任意の可換中心係数環$R$への特殊化は、
分母が原点で非零のスカラー多項式しか含まないため許される。
この最後の降下には全次数のPBW定理を仮定する必要すらない。
ただし周囲の代数のAS正則性や中心元の正則性を
結論したい場合には、別途命題\ref{prop:pg-PBW}などが必要である。
式の剰余は$\hbar=0$を代入すれば直ちに$F$となる。
\end{proof}

\begin{corollary}[対角二次混合を含む全Fermat型]
\label{cor:pg-fermat-all-mixed}
$u$を任意個の中心次数一変数とし、
\begin{equation}
 F=\frac{x^3+y^3+z^3}{3}
 +p\,x^2+q\,y^2+r\,z^2+
 Axy+Byz+Czx+Dx+Ey+Hz+G
 \label{eq:pg-fermat-all-mixed}
\end{equation}
とする。ただし$p,q,r,A,B,C$は中心一次式、
$D,E,H$は中心二次式、$G$は中心三次式である。
$X=x+p$, $Y=y+q$, $Z=z+r$を置き、
\begin{equation}
 \begin{aligned}
 D'&=D-p^2-Aq-Cr,&
 E'&=E-q^2-Ap-Br,&
 H'&=H-r^2-Bq-Cp,\\
 G'&=G+\tfrac23(p^3+q^3+r^3)
     +Apq+Bqr+Crp-Dp-Eq-Hr
 \end{aligned}
 \label{eq:pg-fermat-translation}
\end{equation}
とすると、
$F$は$(X,Y,Z)$に関して
\eqref{eq:pg-fermat-mixed-F}の形となり、
その係数は$A,B,C,D',E',H',G'$である。
従って定理\ref{thm:pg-fermat-mixed-universal}の公式へ
これらを代入し、$X,Y,Z$を$x+p,y+q,z+r$に戻すことで、
\emph{任意の}中心変数混合項を含むFermat型の
明示的な有限多項式中心三次元を得る。
\end{corollary}
\begin{proof}
$F(x,y,z,u)$に$x=X-p$, $y=Y-q$, $z=Z-r$を代入して
三次式を展開すると\eqref{eq:pg-fermat-translation}を得る。
$p,q,r$は中心なので、線形平行移動は交換子を保ち、
Weyl対称化と可換である。
ゆえに定理の中心元をこの座標同型で引き戻せる。
\end{proof}

\begin{remark}[Hessianジェットとの一致と新規性]
\label{rem:pg-fermat-mixed-hessian}
\eqref{eq:pg-fermat-mixed-F}のactive Hessian行列式は
\begin{equation}
 \det\operatorname{Hess}_{x,y,z}(F)
 =2(4xyz-A^2z-B^2x-C^2y+ABC).
\end{equation}
従って公式\eqref{eq:pg-fermat-mixed-Omega}は、
Hessianを用いて次の短い\emph{厳密な}表示にも書き直せる：
\begin{equation}
\boxed{\quad
\Omega_{F,\hbar}
 =\sym\!\left(
 F-\frac{\hbar^2}{12}\det\operatorname{Hess}_{x,y,z}(F)
 +\frac{\hbar^2}{6}ABC
 +\frac{\hbar^4}{18}
 (Az^2+Bx^2+Cy^2-ACx-ABy-BCz)
 \right).
\quad}
\label{eq:pg-fermat-mixed-hessian}
\end{equation}
ここで$\sym$は全active単項式へWeyl対称化を施し、
中央係数はそのままとする。
この式は命題\ref{prop:pg-universal-casimir-jet}の二次Hessian補正と整合し、
残りの純中心項$\hbar^2ABC/6$はCasimir規格化の不定性を表す。
重要なのは、$\hbar^4$補正も含む本式が
有限次ジェットでなく\emph{厳密な中心元}であることである。

式\eqref{eq:pg-fermat-mixed-Omega}自体は
論文Bの関連研究に既に現れたものであり、
今回新たに得た成果は、全ての中心係数を独立な記号として保持した
\emph{普遍的・厳密な中心性の計算証明}と
\eqref{eq:pg-fermat-all-mixed}への明示的拡張である。
先行文献に同じ公式が含まれるかは未確認であり、
無条件な新規性は主張しない。
一般のHesseパラメータ$t$については
定理\ref{thm:pg-universal-hesse-mixed}で
全中心混合方向の短い$\hbar$の多項式公式を既に証明した。
さらに三変数制限が特異または零の場合も、
全体の三次式が滑らかなら
定理\ref{thm:pg-universal-determinant}の
代数的Cramer中心元が存在する。
本節の式はその$t=0$の場合の
簡明な閉公式として位置づける。
\end{remark}
